\chapter{We Often Say More Than Our Words}
\section{A Message: ``You There?''}
First pick up something utterly ordinary: your phone.

The lock screen lights up. A message appears, just two words and a question mark: ``You there?''

It is from your desk mate. You stare, finger hovering above the screen, suddenly unsure how to reply. Say ``Yes,'' and the next message may well be ``Lend me your homework so I can check mine.'' Pretend not to see it, and tomorrow's meeting may be awkward. You notice something odd. This is plainly a question, literally requiring only ``yes'' or ``no.'' Yet what races through your mind is not ``Am I there?'' but ``What does he want?'', ``Should I agree?'', and ``How can I reply without refusing or getting trapped?''

You know every word; the grammar could hardly be simpler. But what the message is doing far exceeds its literal content. It knocks on a door, probes, prepares the way for a request. You read not a sentence but an entire plan that has not yet been voiced.

Think of other examples. Mom texts ``What time is it?'' She probably knows the time; she is asking ``Why aren't you home yet?'' A friend replies in a group chat with nothing but ``Oh.'' You immediately feel something is wrong, though any dictionary gives only an acknowledgment.

That is this chapter's subject: we often say more than our words literally mean. Speech is an iceberg. Everyone can consult a dictionary for the little part above water; what decides whether communication succeeds lies mostly below: hints, presuppositions, irony, requests, promises, and who is entitled to say such things to whom. Let us dive beneath the surface.

\section{One Sentence in a Winter Living Room}
Now come into a scene.

Time: an evening in January, 7:40 p.m. Place: an ordinary family's living room in a northern Chinese city. Wind blows outside; an old aluminum window frame will not close tightly. Air squeezes through the cracks with a faint moan.

Four people are around. Grandma knits on the sofa nearest the window. Mom, just home from work, checks bills on her phone at the dining table. Fifteen-year-old Xiaozhou sits cross-legged on the rug playing a game. Dad washes dishes in the kitchen amid rushing water.

Grandma's hands pause. She looks toward the window and says:

``It's a little cold in here.''

A great deal happens in the next two seconds, though nobody moves much.

Xiaozhou's fingers and eyes stay on his screen. ``I'm not cold,'' he says. He hears a statement about temperature and, as the family's most scientifically informed member, takes responsibility for correcting Grandma's inaccurate report.

Mom does not look up. ``Xiaozhou, go close the window properly.'' She hears a request, translates it into an instruction, and forwards it to the household's lowest-ranking person.

Dad pokes his head out of the kitchen. ``Mom, why don't you move over a little? My hands are full with the dishes.'' He too hears a request, cannot fulfill it, and offers an alternative.

Grandma smiles. ``It's all right. I was only saying.'' She starts retracting the sentence because it has caused more commotion than she intended.

One short sentence, in one room at one instant, becomes three different things: a weather report, a request to close a window, and a household matter requiring someone to be assigned. Its speaker finally declares it none of these, merely ``something I said.''

Remember this living room. Most of the rest of the chapter answers the questions hidden in this little scene.

\section{What Lies Beyond the Literal?}
First lay out the conflict, without rushing to an answer.

Xiaozhou has a case that sounds quite upright: a sentence means what its words say. ``Cold'' means cold; ``here'' means here. Grandma never said ``Please close the window,'' so it is not a request. If everyone must guess what others imply, what rules can speech have? Who is responsible for a wrong guess? He sees himself not as uncaring but honest: I answer only for what is explicitly stated.

Mom's case is just as strong. An ordinary person sitting beside a drafty window tells the family ``It's a little cold in here.'' If that does not suggest wanting the window closed, humans might as well stop expressing needs through language and press buttons instead. Understanding speech has always required more than ears: the occasion, tone, speaker, and reason for speaking just now.

Behind these positions is a real question, not a contest over emotional intelligence:

\textbf{Where does a sentence's meaning live?}

In words and grammar, like goods in a box that are the same whoever opens it? In the speaker's head, with the listener merely unpacking the parcel? Or in the space between two people, assembled for the moment from circumstances, relationships, and shared understanding?

If the first, ``It's a little cold in here'' is always just a weather report. Xiaozhou is right, Mom overinterprets, and listeners owe apologies for every misunderstanding caused by ``I was only saying.'' If the second, meaning can never be checked: the speaker can always say ``That's not what I meant,'' leaving the listener perpetually wrong. If the third, matters get harder still. Meaning needs tacit understanding; where does that come from? Are those who lack it---strangers, newcomers, foreigners, straightforward personalities---inherently disadvantaged?

This is no classroom game of rhetoric. It decides practical questions: whether ``sorry'' counts as an apology, whether a joke is evidence of offense, whether a teacher's ``Think again'' is advice or warning, whether law recognizes an understanding both parties assumed but never wrote into their contract.

Before reading on, take a side in the living room: did Xiaozhou have an obligation to understand Grandma?

\section{Who Speaks, and Who Is Entitled to Speak?}
First give both living-room positions their strongest case, then look at how other places handle the same difficulty.

\textbf{Position one: people should say things directly.}

Let us state Xiaozhou's case fully. ``Direct speech'' is often equated with ``low emotional intelligence'' in our surroundings, which is unfair.

The directness camp has at least three reasons. First, clarity is kindness. State your request plainly, and you return the whole choice to the other person: accept or refuse, without agonizing over ``Is that really what she means?'' Second, directness is fair to people unfamiliar with local rules. Implication depends on shared understanding, a local resource. New classmates, colleagues from elsewhere, and people who find guessing intentions difficult always lose in a ``you work it out'' culture. Explicit speech lowers communication's admission price to zero. Third, directness clarifies responsibility. I say ``Please close the window'' and you do not: responsibility is yours. I say only ``A little cold'' and you do nothing: what right have I to be angry? Much resentment in families and offices comes precisely from ``I thought you understood.''

This is a strong case. In many settings requiring intense cooperation---cockpits, operating rooms, construction sites---people protect lives by eliminating implication. A pilot reports ``Low fuel''; the control tower does not answer ``Oh, do you consider that serious?'' There is no room for hints there.

\textbf{Position two: strictly literal speech leaves many things unsayable.}

Mom's side has a complete case too. Many human needs are soft: wanting company, feeling overlooked, hoping another person will take initiative. State them outright, and their nature changes. ``Please care about me spontaneously'' destroys what it seeks, because solicited care is no longer spontaneous. Indirection is not hypocrisy; it leaves both sides a way out. Grandma says ``A little cold,'' and if nobody responds she can gracefully act as though she said nothing, instead of becoming a rejected beggar. Say directly ``Xiaozhou, close the window for me,'' and being rebuffed becomes an official conflict. Implication is society's shock absorber: it makes refusal cheaper and testing the waters safer.

Look elsewhere and you find this is neither uniquely Chinese nor captured by crude claims that ``Easterners are indirect, Westerners direct.''

Japan has an expression almost everyone uses: ``read the air,'' kūki o yomu. It means reading unspoken expectations from a situation's atmosphere. No objections to your proposal at a meeting may not mean no opposition; people may be waiting for you to ``read'' it yourself. Japanese society values this ability highly; people unable to read the air may be excluded just like those who speak out of turn.

On Java, Indonesia's most populous large island, linguistic hierarchy is built directly into grammar. Javanese has several levels according to the addressee's status. Speaking to a servant and to an elder requires changing vocabulary and sentence patterns together. There, ``literal meaning'' almost ceases to exist as a separate concept: your first word already declares ``who I think you are, and who I am.''

English-speaking cultures are equally full of bends. A British speaker's ``Could you pass the salt?'' literally asks about ability, yet everyone recognizes a request. An American academic writes ``This is interesting'' in comments on a paper; insiders know it often politely means ``I disagree but don't want an argument.'' English-teacher training materials even specifically teach international students not to answer ``Yes, I can pass the salt'' and remain seated.

Then there is familiar Chinese ``Have you eaten?'' As a greeting, it defeats a dictionary alone. Literally a question about food, it functions roughly as ``Hello.'' A conscientious foreigner who stops to report exactly what they ate is not foolish; they possess a dictionary but not the shared understanding.

Consider an ancient Chinese scene. The Zhao section of the Strategies of the Warring States records a famous story. Qin attacks Zhao; Zhao asks Qi for help. Qi demands that the Queen Dowager of Zhao send her younger son, the Lord of Chang'an, as a hostage before troops are dispatched. Furious, she announces that she will spit in the face of anyone mentioning it again. Every minister's remonstrance hits this literal wall. Elder statesman Chu Long enters the palace and never mentions the hostage. He first asks about meals and daily life, then seeks a position for his own younger son, and finally discusses what it means for parents truly to love their children. The dowager herself supplies the crucial sentence: ``Parents who love their children plan far ahead for them.'' By then, without mentioning the hostage again, the matter is settled: she sends the Lord of Chang'an to Qi. Chu Long never violates her literal prohibition, yet achieves what every direct remonstrator failed to do. He has not tricked her with pretty words. He layers the real request beneath preparation, letting her reach the conclusion herself. People more readily accept conclusions they ``work out'' than ones forced upon them.

Thus ``hear what lies behind the words'' is not one culture's quirk but a universal human apparatus. Cultures differ in how deeply it is installed, where it is used, and how costly failure to read it is.

Return to the living room. Notice a detail we passed over: after Grandma says ``It's a little cold,'' who actually closes the window? Probably Xiaozhou. The sentence leaves Grandma, passes through Mom's translation, and lands on the family's youngest, lowest-ranking member. Implication contains power as well as meaning. Some utterances work only on particular people; some people may not utter them. Identical words from different mouths become different things. We will return to this in greater depth.

\section[Thinking Bridge: Layers of an Utterance]{Thinking Bridge: Three Layers of an Utterance}
Can we analyze ``It's a little cold in here'' with a portable tool instead of relying on talent and intuition? Yes. Linguists have long done so. The tool's core is to separate any utterance into three layers.

\textbf{Layer one: what is literally said (semantics).}

Semantics studies words' and sentences' ``standard equipment.'' Cold indicates low temperature, here the speaker's location; the sentence states a fact about felt conditions. Dictionaries and grammar books govern this layer, in principle the same for anyone consulting them. This is what Xiaozhou hears.

\textbf{Layer two: what is taken for granted (presupposition).}

A presupposition is background beneath the sentence that the speaker treats as true without arguing for it. Grandma's ``It's a little cold in here'' already assumes that it could be warmer, that cold is not normal, and that someone present can change it. A sharper example: ``Have you finished your homework?'' presupposes that you have homework. If none was assigned, the answer is not merely negative; the whole question collapses. Presuppositions are dangerous because simple denial leaves them intact. ``I haven't finished'' preserves the assumption. Only ``Wait, there was no homework today'' dismantles it. Advertisers and interrogators understand this: answer either yes or no to ``Are you still staying up late gaming?'' and you have already conceded that you did so before.

\textbf{Layer three: what the utterance is doing (pragmatics).}

Pragmatics asks what this sentence, spoken here, by this person to that person, actually accomplishes. Statement, request, complaint, probe? In this living room, addressed to these family members, ``It's a little cold'' requests that the window be closed. This is Mom's layer.

In the mid-twentieth century, philosopher Paul Grice proposed a useful idea: \textbf{conversational cooperation}. He suggested that speakers tacitly sign an invisible contract with roughly four understandings: give enough information (quantity); say what you believe is supported (quality); stay relevant (relation); and be as clear as possible (manner). The ingenious part is that implications are conveyed precisely \textbf{by departing from these understandings}. Grandma plainly knows whether she is cold. Said at this moment, her sentence looks suspicious in information and relevance. Because the literal reading seems insufficient, listeners are pushed to infer what she wants: close the window. Hints are not exceptions to cooperation but its engine.

Irony drives the same engine harder. You fail an exam, and a friend says ``You're really something.'' Literally praise, even extravagant praise, sits beside plainly poor results and sharply breaches ``say what you believe.'' Meaning flips. Irony generates power from the drop between two layers.

Next time a sentence makes your heart sink, do not rush into anger or hurt. Try the three-layer analysis: what does it say, what does it assume, what does it do? Unpack many arguments, and you find two people living on different floors.

\section{Two Millennia of Distrusting Fine Words}
Now read a little primary material. More than two thousand years ago, people already worried about what lies beyond the literal.

The Analects, in ``Yang Huo,'' records Confucius saying:

\begin{quote}
The Master said: ``Clever words and a pleasing countenance seldom accompany benevolence.''
\end{quote}
Roughly: people with polished speech and carefully pleasant faces rarely possess genuine moral goodness. The Analects voices similar cautions elsewhere. ``Zi Lu'' contains an even more famous passage:

\begin{quote}
``If names are not correct, speech does not accord; if speech does not accord, affairs cannot succeed.''
\end{quote}
In essence: unless titles and names are set right, speech cannot follow properly; without proper speech, things cannot be accomplished.

Read the passages together and notice the tension. The first distrusts language's surface: literal words and expressions can be packaging; the prettier the wrapper, the more carefully you should inspect the goods. The second values that surface intensely: even apparently formal matters like ``what to call something'' affect success. Wrong names disrupt speech, which disrupts action.

Are they contradictory? Not necessarily. Together they describe two sides of one matter. Confucius doubts that language automatically equals inner feeling, hence distrust of clever words. Yet he believes language's use can shape reality, hence rectifying names. A title or expression is not merely a label describing the world; it changes the world. Children describe parents' mistakes as ``for your own good''; a court describes conquest as ``comforting the people and punishing the guilty.'' Once a name changes, so do the act's apparent nature and whether bystanders should be angry.

This leads to a question still making headlines two millennia later: is language reality's mirror or its construction crew? Confucius seems to answer: it pretends to be a mirror while frequently doing secret construction. His defense is not silence, but attention to the gap between words and inner motives, names and actual conditions. The wider the gap, the greater the caution needed.

We use the Analects as evidence not because Confucius must be right, but to show that your group-chat anxiety over ``heh heh'' belongs to one of humanity's oldest worries. It has a documented history and was ethical as well as linguistic from the beginning. We will see that it is also about power.

\section{One Sentence, Three Explanations}
Return to ``It's a little cold in here.'' We now have more parts with which to build different explanations of the same fact. Distinguish four things: what happened (Grandma's short sentence and Xiaozhou's ``I'm not cold''), the relatively undisputed evidence; why it happened, where explanations compete; how participants understood it (Grandma says she was only speaking casually), their own account; and how we evaluate it (was Xiaozhou inconsiderate?), a value judgment. We compare the second: why the scene unfolded.

\textbf{Explanation one: inference went wrong (the Gricean account).}

On the cooperative view, Grandma offers a standard hint. Its questionable literal relevance guides listeners toward ``Please close the window.'' Xiaozhou misses it because his inference chain breaks: either he never starts, absorbed in his game, or deliberately cuts it short, unwilling to close the window. This account is clear and testable, turning ``understanding'' from mystery into a reasoning process checked step by step. Its limit is clear too. It assumes two people calmly solving a puzzle, yet this living room's occupants are not equally situated. Why does Grandma not simply say ``Xiaozhou, close the window''? The account explains how to infer indirect speech, not why people choose it.

\textbf{Explanation two: face is at work (politeness theory).}

Anthropologists and linguists note two conflicting desires almost everyone has: to be liked and not rejected, ``positive face,'' and to avoid imposition and retain freedom, ``negative face.'' A direct command can damage both: the commander seems overbearing, the commanded unfree. Hints are face's adhesive bandage. The vaguer the request, the cheaper refusal becomes and the safer both faces remain. Here Grandma is not unable to speak directly. The family jointly protects everyone's face; Xiaozhou's ``I'm not cold'' is the real disruption, refusing this gentle collaboration. This fills the prior account's gap: why humans \textbf{need} implication. But it has its own blind spot. Treating everyone as a shrewd individual face-manager cannot explain the unequal labor involved. Why must Grandma circle around her needs while Xiaozhou can directly say ``I'm not cold''?

\textbf{Explanation three: power determines sentence form (the social-structure account).}

This account changes viewpoints. It asks not ``How is this understood?'' but ``Who is permitted to speak how?'' Mom may say ``Xiaozhou, close the window'' because she has authority over the child. Grandma uses tentative language with her son and daughter-in-law, especially about changing things in her son's home, because she lives there as a guest and has a reduced entitlement to ask. Xiaozhou dares answer directly because he is the indulged center. Sentence forms map power: higher status can afford imperatives; lower status disguises requests as statements. Elsewhere it is similar. An employee says ``This proposal may still have room for improvement''; the boss can say ``Redo it.'' The employee cannot say that to the boss. This sharp explanation captures the asymmetry the others miss. Its limit: reduce everything to power and we overlook genuine kindness. Grandma may simply feel cold while caring about the child; her indirectness may contain love, not only subordination. Translating family feeling entirely into power calculations can be another misreading.

Each account holds part of the truth and fails to reach all of it. This is not fence-sitting but a methodological reminder: an explanation that explains everything has often discarded something in advance.

\section{Further Challenge: Speaking Is Acting}
Now introduce the chapter's weightiest theory, from British philosopher J. L. Austin. His Harvard lectures in the 1950s later became a book whose title is itself a declaration: How to Do Things with Words.

Austin focused on a peculiar phenomenon philosophers had long neglected. They assumed sentences primarily describe the world, making descriptions true or false: ``The window is drafty.'' But consider these:

\begin{itemize}
\item ``I promise to return your book tomorrow.''
\item ``I declare the competition open.''
\item ``I apologize for what I said.''
\item ``I name this ship Endeavour.''
\item A judge says: ``I sentence you to three years in prison.''
\end{itemize}
Are these true or false? Someone saying ``I promise'' does not \textbf{report} an inner commitment as a thermometer reports temperature. In uttering the words, they \textbf{perform} the promise \textbf{at that moment}. Once spoken, the world changes: you owe a book, the competition clock starts, the ship has a name, someone loses three years of freedom. Austin called this \textbf{doing things with words}: speaking is not merely speaking, but action in itself.

This immediately explains many stubborn facts of life. Why is apology difficult? ``I'm sorry'' is not a report of pain but a procedure: it officially changes the accounts between you and another person, acknowledging a debt. People who feel the hurt but cannot speak are blocked not in description but in daring to \textbf{act}. Why does everyone hold their breath for ``I do'' when an officiant asks at a wedding? Because this is not an information request but a contract signed live.

Austin's most interesting cases, however, are \textbf{speech acts that misfire}. Doing things with words requires conditions. If they are absent, the words accomplish nothing: only sound remains. Consider examples.

A random passerby grabs the wedding microphone and shouts ``I pronounce you married.'' Everyone laughs, but nobody thinks a marriage has thereby occurred. Why? This declaration needs an authorized speaker, such as a registrar or officiant; an occasion, the ceremony; and a procedure, including both parties' assent. Without the right authority, setting, or procedure, the same sentence is merely sound waves.

Consider an easily misjudged counterexample. You might infer that all social consequences therefore depend on who said what, and we need only hear correctly. But imagine a teacher looking at a dozing Xiaoming and saying ``Everyone seems to have rested well last night.'' The class hears the irony; Xiaoming blushes. Literal analysis and pragmatic inference succeed, and the speech act, an indirect criticism, takes effect. Yet is it \textbf{justified}? Before the class, the teacher exercises a power only the teacher enjoys: public mockery without being mocked back. Understanding the implication does not settle whether the act is acceptable. Pragmatics tells you what an utterance accomplished; it cannot make the ethical judgment for you. A tool is a map, not a verdict.

Speech-act theory also reframes identity. We usually think identity comes first, speech afterward: you may sentence because you are a judge. Austin reveals flow in the other direction. Identity is also \textbf{spoken into being}. Only after ``I declare you elected'' takes effect do you become chair. Every class in which only the teacher, not the students, may say ``Quiet'' speaks and reinforces the teacher's identity again. Language and power are not separate things: language is power's most frequently worn clothing.

Weigh this theory's strength and limits. It explains promises, naming, verdicts, apologies, and commands. Does it explain consolation? Two people lying together watching clouds and chatting? Not all speech conducts business; some simply keeps company. Theory draws a bright line, clear on one side. The other---where speech can be only warmth---it leaves to literature.

\section{Literal Battles in the Group-Chat Age}
This question, more than two thousand years old, explodes daily in your pocket.

Online communication magnifies every problem in the chapter. Text chat strips away tone, expression, and pauses, leaving the literal, precisely the least dependable layer. The results follow.

A period becomes a weapon. In countless chats, ``Okay'' and ``Okay.'' signal different attitudes. Young people know that period's weight, though no grammar book says ``a period indicates coldness.'' Users invent a pragmatic layer: when the system supplies only literal words, punctuation, message length, and response speed are recruited for new implications. Fast or slow, one word or a paragraph, ``mm'' or ``mm-hmm'': each is a code, and the codebook keeps changing.

Irony runs uncovered. Online snideness institutionalizes irony: praise on the surface, disparagement beneath, literal compliments delivering attacks. It thrives precisely because platform rules detect literal content. Automated moderation can remove swearwords but not ``Aren't you a genius.'' People invent words; words evade machines; machines drive people to invent more. This is an arms race between the literal and the pragmatic.

Presupposition becomes standard equipment. ``Why are young people today all unwilling to work hard?'' Any platform can generate thousands of responses, yet few first dismantle the assumptions: were young people formerly willing to suffer? Has the relationship between suffering and success been settled? Today's version of the ``incorrect names'' Confucius feared is that a trending topic's wording has already cast a secret vote. Half the battle is decided before discussion starts. Asking ``What does this question take for granted?'' is basic self-defense in the information flood.

Speech acts are also mass-produced online, often with corners cut. Recall institutional or celebrity ``apologies'': ``We deeply regret any inconvenience caused.'' Our tools reveal the mechanisms. Who caused it? The subject disappears; nobody admits ``I did it.'' ``Inconvenience,'' not ``wrongdoing,'' acknowledges your discomfort rather than my fault. ``Deep regret'' is a hollowed-out speech act, performing apology's procedure while removing its core: owning the debt. Austin taught that speech acts require conditions. Skilled statement writers know this and extract those conditions one by one, leaving an apology-shaped shell. You can now see through it. That is what a tool offers.

The newest variable is that your conversation partner may not be human at all. With a chatbot, the whole pragmatic apparatus can fail. It does not read hesitation in your ``You there?''; its ``I understand how you feel'' accompanies no understanding. It is a promise without commitment, an automatically performed speech act. This does not mean every machine utterance is false. It means machines force us into a strange position: for the first time, we see with such clarity how much we ordinarily seek in human speech that is absent from the words themselves.

\section{Over to You: Literal Speech Only?}
Return to the opening living room.

On the evening Grandma says ``It's a little cold in here,'' suppose the household has an iron rule: everyone may speak only literally; requests and refusals must be explicit; irony and hints are void. Would this be a better home?

Give your answer and reasons. Before concluding, weigh these considerations once more.

You have directness's case: clarity, fairness, clear responsibility, protection for newcomers, and the life-preserving communication of cockpits and operating rooms.

You also have indirection's case: a graceful retreat from refusal, dignity for needs, room for soft and embarrassing feelings, and ways for lower-status people to avoid head-on collisions every time.

You know the harder fact too: power hides in implication. Who circles around and who speaks directly often depends on status, not personality. Yet translating everything into power overlooks the real warmth in Grandma's words: simply feeling cold while caring about the child.

So is a world allowing only direct speech more honest or more cruel? Does it lower the vulnerable person's admission price to zero, or remove their only shock absorber? Is indirectness consideration, or selfishly transferring the cost of communication to the listener?

There is no standard answer. But notice that every sentence you use to answer is doing something beyond the literal: taking sides, probing, asking to be understood. That has been happening since you opened the first chapter. Starting today, you have eyes to see it.
